\section{Schéma électrique}
\label{sec:schema}

Le schéma électrique réunit les composants retenus à la section précédente et fixe leurs interconnexions, des étages de puissance jusqu'aux protections. Cette section en décrit d'abord l'organisation, puis les sous-circuits qui demandent une explication propre, protections comprises. Elle ne reprend ni la justification des composants, traitée à la section~\ref{sec:choix_composants}, ni le routage, abordé à la section~\ref{sec:layout}. Les feuilles complètes du schéma figurent en annexe~\ref{ann:schema}.


\subsection{Organisation du schéma}

Le schéma est hiérarchique et se lit de l'extérieur vers l'intérieur. La première feuille rassemble les connecteurs accessibles depuis l'extérieur de la carte, avec les protections qui leur sont associées, et reprend ainsi le découpage de l'analyse fonctionnelle de niveau 0 présentée à la figure~\ref{Analyse_Fonctionnelle-Niv1.drawio}. La feuille suivante descend d'un niveau et expose les six blocs principaux du système, à savoir l'alimentation, le microcontrôleur, l'écran, l'interface à boutons, le pont USB-UART et la chaîne audio filaire, dans la continuité de l'analyse fonctionnelle de niveau 1 de la figure~\ref{Analyse_Fonctionnelle-Niv2.drawio}. Chaque bloc renvoie à sa feuille propre, ce qui suit les frontières naturelles du système et permet de vérifier chaque sous-ensemble séparéement.

\subsection{Distribution de l'alimentation et chaîne de charge}

L'énergie entre par le connecteur USB-C et rejoint directement le chargeur MAX77757, qui détermine seul le courant maximal qu'il peut tirer de la source à laquelle il est relié. Il dispose pour cela de deux mécanismes de détection complémentaires, l'un s'appuyant sur les lignes de configuration du connecteur pour reconnaître une source Type-C et le courant qu'elle annonce, l'autre sur les lignes de données pour identifier un chargeur mural classique. Les résistances de terminaison exigées par le premier sont intégrées au chargeur et n'apparaissent donc pas sur le schéma~\cite{analog_MAX77757}. Les lignes de données, en revanche, ne lui parviennent pas directement et passent d'abord par un multiplexeur analogique.

Ce multiplexeur répond à une contrainte de partage. Deux composants ont besoin de la même paire D+ et D$-$, le pont USB-UART pour la programmation et le chargeur pour sa détection. Mais une paire USB ne se partage pas entre deux périphériques sans dégrader l'impédance de ligne et sans risquer que les deux répondent en même temps. Le microcontrôleur ne peut pas non plus arbitrer lui-même, l'ESP32 ne disposant d'aucun contrôleur USB natif. Un multiplexeur analogique PI3USB221A aiguille donc la paire vers l'une ou l'autre destination selon un signal de commande issu du microcontrôleur. Au repos, une résistance de pull-down relie les données USB au pont USB-UART, ce qui permet de programmer la puce et d'accéder à la console série dès la mise sous tension. Sur commande, il bascule la paire vers le chargeur, qui peut alors mener sa détection BC1.2 afin de reconnaître un chargeur mural. Un seul port physique dessert ainsi tour à tour la programmation et la négociation de charge, sans jamais connecter les deux fonctions en même temps~\cite{diodes_inc_PI3USB221A}.

Le MAX77757 gère la distribution de l'énergie entre l'entrée, le système et la batterie grâce à son étage power path intégré. Ce module priorise en permanence l'alimentation du système et réserve à la charge de la batterie uniquement la puissance disponible sous la limite de courant négociée. Cette architecture garantit un fonctionnement continu même avec une batterie déchargée et assure une transition invisible vers la batterie lors du retrait de la source externe.

La jauge MAX17260 s'insère sur ce chemin pour suivre l'état de la batterie. Elle combine deux grandeurs, la tension de la cellule et la charge échangée, mesuré aux bornes d'une résistance de mesure placée en série avec la batterie. Cette résistance est placée du côté haut, entre le pôle positif de la cellule et le nœud système, plutôt que dans le retour de masse. Le plan de masse reste ainsi continu et unique, alors qu'une mesure côté bas insérerait la résistance dans le chemin de retour et créerait entre la masse de la cellule et celle de la carte un écart de potentiel proportionnel au courant consommé, dégradant la référence de la chaîne audio. Les grandeurs mesurées et l'état de charge estimé sont transmis au microcontrôleur par I\textsuperscript{2}C, et une broche d'alerte signale au firmware le passage sous un seuil bas de charge.

\fig[H, width=0.9\linewidth]{Mesure de courant côté haut, position de la résistance de mesure entre la cellule et le nœud système. Diagramme créé avec l'aide de Claude Opus 4.8.}{mesure_courant.drawio}

Le rail de 3,3~V est produit à partir de la tension \texttt{VSys} par le convertisseur abaisseur TPS62A01. Les composants passifs qui l'entourent ont été calculés d'après sa fiche technique pour une tension de sortie de 3,3 V, à savoir le pont de retour qui fixe la consigne à partir de la tension de référence interne, et le condensateur d'avance de phase placé en parallèle sur sa résistance haute, qui réduit l'ondulation résiduelle en mode d'économie d'énergie et améliore la réponse aux échelons de charge~\cite{TI_TPS62A01}. Comme exposé à la section~\ref{sec:topologie_buck}, en fin de décharge le convertisseur se comporte en régulateur à très faible chute, ce qui maintient le bas de la plage batterie exploitable.

La figure \ref{arbre_alim.drawio} montre l'arbre d'alimentation du circuit.

\fig[H, width=0.9\linewidth]{Arbre d'alimentation, de l'entrée USB au rail logique. Diagramme créé avec l'aide de Claude Opus 4.8.}{arbre_alim.drawio}

L'entrée d'activation du convertisseur est pilotée par une logique câblée, détaillée à la figure~\ref{extrain_enable.png}, qui combine trois sources de mise en marche validées par le superviseur de tension.

\fig[H, width=1\linewidth]{Extrait du schéma, logique d'activation du convertisseur et superviseur de tension}{extrain_enable.png}


\subsection{Chaîne audio filaire}

La chaîne audio filaire conduit le flux numérique du microcontrôleur jusqu'à la prise casque. Sa partie analogique, de la sortie différentielle du convertisseur au jack, est reprise sous forme simplifiée à la figure~\ref{chaine_audio.drawio}, la feuille complète du schéma figurant à l'annexe~\ref{ann:schema}.

\fig[H, width=1\linewidth]{Schéma de principe simplifié d'un seul canal de la chaîne analogique, de la sortie différentielle du convertisseur au jack, le second canal étant strictement identique. Diagramme créé avec l'aide de Claude Opus 4.8.}{chaine_audio.drawio}

\subsubsection{Liaison numérique}

Le microcontrôleur attaque le convertisseur par une liaison I\textsuperscript{2}S à trois fils, l'horloge de bit, l'horloge de mot et les données. Aucune piste d'horloge maître n'a été routée entre les deux composants, le PCM5242 reconstruisant sa propre horloge interne à partir de l'horloge de bit au moyen de sa boucle à verrouillage de phase~\cite{TI_PCM5242}. Ce choix retire une piste numérique à haute fréquence de plus sur une carte dense et supprime une source de rayonnement, au prix d'une fréquence d'échantillonnage minimale exploitable, plus de détails sont donnés à la section \ref{sec:horloges_dac}.

\subsubsection{Filtre de sortie}

Chaque ligne de sortie du convertisseur porte deux résistances série de 5,1~k$\Omega$ encadrant un condensateur de 100~pF monté entre les deux lignes différentielles. Cet ensemble forme un filtre passe-bas du premier ordre dont la fréquence de coupure différentielle se situe autour de 312~kHz.

Le nœud du condensateur n'est pas un point à vide. La seconde résistance de chaque ligne rejoint l'entrée inverseuse de l'amplificateur, que la contre-réaction maintient au potentiel de référence et qui se comporte donc en masse pour le signal. Vue de ce nœud, la résistance amont apparaît ainsi en parallèle avec la résistance aval, ce qui ramène la résistance de chaque branche à $R \parallel R = 2{,}55\,\text{k}\Omega$. Le condensateur étant monté entre les deux lignes, le signal différentiel traverse cette résistance deux fois. La fréquence de coupure vaut donc
\begin{equation}
  f_c = \frac{1}{2\pi \cdot 2\,(R \parallel R) \cdot C} = \frac{1}{2\pi \times 5{,}1\,\text{k}\Omega \times 100\,\text{pF}} \approx 312\,\text{kHz},
\end{equation}

Un ordre aussi faible suffit parce que le PCM5242 est un convertisseur delta-sigma à suréchantillonnage. Son filtre d'interpolation numérique porte la cadence de sortie à 352,8 ou 384~kHz pour les familles de fréquences usuelles~\cite{TI_PCM5242}. La première image d'un signal de 20~kHz se trouve ainsi reportée à 332,8~kHz, alors qu'un convertisseur sans suréchantillonnage cadencé à 44,1~kHz la placerait à 24,1~kHz et exigerait un filtre analogique raide pour l'éliminer. Le rejet des images est donc entièrement acquis avant l'étage analogique.

Ce qui reste à la sortie du DAC est le bruit de quantification mis en forme par le modulateur, dont l'énergie augmente avec la fréquence et que la fiche technique caractérise de 100~kHz à 3~MHz. La fiche technique indique que de nombreuses applications exigent un filtre RC externe pour en obtenir un rejet suffisant, faute de quoi ce bruit engendre des perturbations électromagnétiques et du repliement dans les étages en série, et signale que ce même filtre améliore la protection des sorties contre les décharges électrostatiques~\cite{TI_PCM5242}. C'est cette tâche que remplit le filtre, sa pente de six décibels par octave donnant une atténuation plus grande sur les décades supérieures.

La fréquence de coupure découle de deux contraintes. Les résistances sont fixées par la charge présentée en sortie du DAC, dont les performances sont caractérisées sur 10~k$\Omega$~\cite{TI_PCM5242}. Des résistances de 5,1~k$\Omega$ étaient donc le choix le plus adapté. Le condensateur, lui, se trouve sur le nœud d'entrée de l'amplificateur, où une valeur élevée réduirait la marge de phase de la boucle de contre-réaction. La valeur de 100~pF du schéma de référence est donc conservée~\cite{TI_PCM5242}.

Cette fréquence de coupure reste cohérente avec le contenu à rejeter. Le bruit gênant se situe dans le domaine du mégahertz, où il produit des perturbations rayonnées et du repliement dans les étages en aval. Un filtre du premier ordre y atténue d'environ 20~dB à 3~MHz, avec une pente qui se poursuit au-delà. Une coupure plus basse renforcerait ce rejet, mais exigerait un condensateur plus grand sur le nœud sommateur ou un filtre du second ordre, deux options écartées par la contrainte de stabilité et par l'encombrement.

L'influence théorique du filtre sur la bande audio reste par ailleurs négligeable. En posant $x = f/f_c$, le filtre présente à 20~kHz
\begin{equation}
  20\log_{10}|H| = -0{,}02\,\text{dB}, \qquad \varphi = -\arctan(x) = -3{,}7^\circ, \qquad \tau_g = \frac{\tau}{1+x^2} = 0{,}51\,\mu\text{s},
\end{equation}
avec $\tau = 0{,}51\,\mu\text{s}$. Le déphasage correspond à un temps de propagation qui ne varie que de 0,4~\% entre le continu et 20~kHz, soit un simple retard identique sur les deux canaux et non une distorsion de phase. Les mêmes résistances constituent enfin la résistance d'entrée qui fixe le gain de l'amplificateur décrit ci-après.

\subsubsection{Adaptation de niveau et choix du gain}

Le PCM5242 délivre une tension différentielle de pleine échelle $V_{\text{DAC}} \approx 4{,}2\,\text{V}_{\text{RMS}}$ à $0\,\text{dBFS}$~\cite{TI_PCM5242}. Or, sous une alimentation de $3{,}3\,\text{V}$, la tension de sortie maximale de l'amplificateur est limitée à $V_{\text{amp,max}} = 2{,}15\,\text{V}_{\text{RMS}}$ sur une charge de $600\,\Omega$~\cite{analog_MAX97220A}. Transmettre directement la pleine échelle du convertisseur reviendrait à demander à l'étage de sortie une tension supérieure à ce qu'il peut reproduire, ce qui impose une adaptation de niveau entre les deux étages.

Le gain de l'amplificateur est fixé par le rapport des résistances externes selon $A_v = \frac{R_f}{R_{in}}$.

La résistance de contre-réaction vaut $R_f = 5{,}1\,\text{k}\Omega$, tandis que la résistance d'entrée est formée des deux résistances série de la ligne, soit $R_{in} = 5{,}1\,\text{k}\Omega + 5{,}1\,\text{k}\Omega = 10{,}2\,\text{k}\Omega$. Ce qui donne :
\begin{equation}
  A_v = \frac{5{,}1\,\text{k}\Omega}{10{,}2\,\text{k}\Omega} = 0{,}5 \quad (\approx -6\,\text{dB}),
\end{equation}
\begin{equation}
  V_{\text{out,max}} = A_v \cdot V_{\text{DAC}} = 0{,}5 \times 4{,}2\,\text{V}_{\text{RMS}} = 2{,}1\,\text{V}_{\text{RMS}} \le V_{\text{amp,max}}.
\end{equation}
La pleine échelle numérique du convertisseur est de la sorte calée juste sous la tension de sortie maximale de l'amplificateur. Cette comparaison porte toutefois sur la charge de référence de 600 $\Omega$ du datasheet. Sur une charge de casque, la contrainte n'est plus une tension mais une puissance, ce que la mesure a mis en évidence à la section~\ref{sec:res_audio}.


\subsection{Programmation et séquence de démarrage}
\label{sec:autoflash}

Cinq broches de l'ESP32 sont des \gls{strapping}, dont le niveau logique au reset fixe le mode de démarrage, avant de redevenir des entrées-sorties ordinaires~\cite{ESP32_WROVER-E_Datasheet}. Le brochage du circuit les utilise toutes, les signaux ont donc été répartis de sorte que leur état de repos coïncide avec la valeur attendue au démarrage, les résistances internes et externes de chaque ligne étant choisis en conséquence. Les deux bus \gls{spi}, actifs à l'état bas, occupent ainsi les pins dont le défaut est haut, tandis que le reset de l'écran et sa ligne C/D occupent celles qui doivent rester basses. Cette dernière contrainte est la plus stricte, GPIO12 sélectionnant la tension d'alimentation de la mémoire flash et de la \gls{psram} du module, un niveau haut au reset rend le démarrage impossible~\cite{ESP32_WROVER-E_Datasheet}.

Le passage automatique en mode programmation s'appuie sur le bootloader de premier niveau intégré à la ROM de l'ESP32. Ce bootloader gère le téléchargement par UART indépendamment du contenu de la mémoire flash et s'active lorsque la broche GPIO0 est maintenue à l'état bas lors du reset~\cite{ESP32_Datasheet}. Le circuit exploite les lignes de contrôle DTR et RTS du CP2102N pour piloter les broches EN et GPIO0 via deux transistors selon le schéma de la figure~\ref{extrait_autoflash.png}. Un réseau RC retarde légèrement l'établissement de GPIO0 par rapport à EN afin de garantir l'ordonnancement correct de la séquence de reset. Cette architecture reproduit celle de la carte de développement officielle d'Espressif~\cite{ESP32_DevKitC_Doc}.

\fig[H, width=0.65\linewidth]{Extrait du schéma, circuit de passage automatique en mode programmation}{extrait_autoflash.png}

\subsection{Conditionnement du potentiomètre de volume}

\begin{wrapfigure}{l}{0.35\linewidth}
  \centering
  \includegraphics[width=1\linewidth]{\assetsdir/potentiometre.png}
  \caption{Montage du potentiomètre de volume, conditionné par deux résistances d'adaptation vers l'entrée ADC}
  \label{fig:potentiometre}
\end{wrapfigure}

Le réglage de volume repose sur un potentiomètre de 10~k$\Omega$ dont le curseur est lu par l'\gls{adc} du microcontrôleur. Cette lecture n'est fidèle que sur une plage restreinte, la fiche technique de l'ESP32 donnant une plage de mesure utile de 150 à 2450~mV avec l'atténuation maximale, la précision se dégradant au-delà~\cite{ESP32_Technical_Reference_Manual}. Relier directement le potentiomètre entre le rail de 3,3~V et la masse placerait donc les deux extrémités de la course hors de cette plage, avec un curseur saturé sur une portion notable de la rotation.

Le montage retenu, donné à la figure~\ref{fig:potentiometre}, encadre pour cette raison le potentiomètre par deux résistances : R11 de 3,3~k$\Omega$ vers le rail et R67 de 620~$\Omega$ vers la masse, formant un diviseur dont le potentiomètre constitue la portion centrale. Un condensateur de 100~nF filtre par ailleurs le bruit capté par la piste du curseur. La tension du curseur est ainsi bornée par les deux valeurs extrêmes suivantes, calculées à partir des résistances du diviseur : \\

\begin{equation}
  V_{\text{min}} = 3{,}3 \times \frac{620}{13\,920} \approx 147~\text{mV},
  \qquad
  V_{\text{max}} = 3{,}3 \times \frac{10\,620}{13\,920} \approx 2{,}518~\text{V}.
\end{equation}
La quasi-totalité de la rotation tombe de la sorte dans la plage exploitable de l'\gls{adc}, seuls les derniers pour cent de la course dépassant légèrement la borne haute, sans conséquence puisque le volume y est déjà à son maximum.

\subsection{Bouton de mise en marche}

Le bouton A remplit deux fonctions, celle d'entrée utilisateur lue par le microcontrôleur et celle de commande de mise en marche du système. Cette seconde fonction interdit de le relier au rail de 3,3~V, qui n'existe pas tant que le régulateur n'a pas démarré. Le bouton est donc raccordé à \texttt{Vsys}, disponible dès qu'une batterie est présente, ce qui lui permet de réveiller la carte depuis l'état éteint, selon le câblage donné à la figure~\ref{bouton_a.png}.

Ce raccordement expose en contrepartie l'entrée du microcontrôleur à une tension supérieure à sa limite, la tension système atteignant 4,2~V sur une batterie pleine. Une diode Zener de 3,3~V placée entre la ligne du bouton et la masse écrête cette tension, la résistance série de 10~k$\Omega$ limitant le courant dans la Zener. Un condensateur de 100~nF et une résistance de rappel de 100~k$\Omega$ complètent le réseau, le premier pour l'antirebond matériel et la seconde pour garantir un niveau bas défini bouton relâché.

\fig[H, width=0.65\linewidth]{Réseau de protection et d'antirebond du bouton A, alimenté par la tension système et écrêté par une diode Zener}{bouton_a.png}


\subsection{Protections électriques}

Comme la carte est exposée aux manipulations, aux câbles externes et aux variations de la source USB, plusieurs protections sont réparties aux points sensibles.

Les cinq lignes exposées du connecteur USB, à savoir VBUS, les deux lignes de données et les deux lignes de configuration, sont chacune protégées par une diode \gls{tvs} qui reste inactive en fonctionnement normal et écrête toute surtension transitoire vers la masse avant qu'elle n'atteigne les circuits~\cite{texas_instruments_system-level_2025}.

La sortie casque reçoit une protection similaire, complétée d'un filtrage. Chaque canal traverse une perle de ferrite associée à un condensateur, ce qui bloque les perturbations radiofréquence susceptibles d'être captées par le câble du casque puis réinjectées dans l'amplificateur, et une diode de protection écrête les décharges appliquées sur la prise~\cite{microchip_technology_inc_an1767_2014}. L'amplificateur retenu intégrant une protection contre les courts-circuits, aucune résistance série n'est ajoutée sur le chemin du signal.

Les contacts de la carte microSD, accessibles à l'utilisateur, sont eux aussi protégés par un réseau de diodes qui écarte les décharges lors de l'insertion.

Les protections liées à l'énergie sont pour l'essentiel assurées de façon autonome par le chargeur, qui surveille la surtension d'entrée et régule la charge en fonction de la température de la cellule lue sur une thermistance. Un fusible réarmable placé en série avec la batterie complète ce dispositif en cas de surintensité. La protection contre la décharge profonde est quant à elle assurée par le superviseur de tension décrit plus haut.

Enfin, chaque broche d'alimentation des circuits actifs reçoit un découplage local, conforme aux schémas de référence des fiches techniques, afin de fournir les appels de courant rapides au plus près des composants et de contenir le bruit de commutation~\cite{prof_bogatin_rules_2013}.


\subsection{Validation du schéma}

Les outils de l'environnement de conception ont d'abord validé le schéma en détectant les broches non connectées, les conflits de sortie et les réseaux isolés. Une vérification manuelle a ensuite confronté chaque connexion aux fiches techniques des composants, avec une attention particulière aux réseaux de configuration résistifs du chargeur et du convertisseur abaisseur dont les valeurs conditionnent un comportement irréversible par logiciel. Enfin, M. Pascal Zosso a procédé à une relecture externe.

Un défaut de conception a été identifié lors des premiers tests sur la ligne du bouton A. Cette ligne pilotant à la fois une entrée du microcontrôleur et l'anode de la diode d'activation du convertisseur abaisseur, l'appui sur le bouton carte éteinte empêchait le démarrage. Le microcontrôleur non alimenté conduit alors sa diode de protection interne vers un rail inactif, rabattant la tension à environ 0,6~V. Cette valeur est trop faible comparé aux 1,2~V requis pour l'activation du convertisseur~\cite{TI_TPS62A01}.

La correction déplace l'anode de cette diode en amont de la résistance série, directement sur la broche commutée de l'interrupteur, où l'impédance quasi nulle du contact fermé impose la tension \texttt{VSys} sans que l'écrêtage du microcontrôleur puisse la retenir. Elle a été réalisée par fil volant sur la carte et devra être reportée au schéma lors d'une éventuelle révision, la version annexée à ce rapport étant celle qui a été fabriquée.

Hormis ce point, le schéma est conforme.